% Options for packages loaded elsewhere
\PassOptionsToPackage{unicode}{hyperref}
\PassOptionsToPackage{hyphens}{url}
\PassOptionsToPackage{dvipsnames,svgnames,x11names}{xcolor}
%
\documentclass[
  11pt,
  a4paper,
]{ltjsarticle}
\usepackage{amsmath,amssymb}
\usepackage{iftex}
\ifPDFTeX
  \usepackage[T1]{fontenc}
  \usepackage[utf8]{inputenc}
  \usepackage{textcomp} % provide euro and other symbols
\else % if luatex or xetex
  \usepackage{unicode-math} % this also loads fontspec
  \defaultfontfeatures{Scale=MatchLowercase}
  \defaultfontfeatures[\rmfamily]{Ligatures=TeX,Scale=1}
\fi
\usepackage{lmodern}
\ifPDFTeX\else
  % xetex/luatex font selection
\fi
% Use upquote if available, for straight quotes in verbatim environments
\IfFileExists{upquote.sty}{\usepackage{upquote}}{}
\IfFileExists{microtype.sty}{% use microtype if available
  \usepackage[]{microtype}
  \UseMicrotypeSet[protrusion]{basicmath} % disable protrusion for tt fonts
}{}
\makeatletter
\@ifundefined{KOMAClassName}{% if non-KOMA class
  \IfFileExists{parskip.sty}{%
    \usepackage{parskip}
  }{% else
    \setlength{\parindent}{0pt}
    \setlength{\parskip}{6pt plus 2pt minus 1pt}}
}{% if KOMA class
  \KOMAoptions{parskip=half}}
\makeatother
\usepackage{xcolor}
\usepackage[margin=24mm]{geometry}
\setlength{\emergencystretch}{3em} % prevent overfull lines
\providecommand{\tightlist}{%
  \setlength{\itemsep}{0pt}\setlength{\parskip}{0pt}}
\setcounter{secnumdepth}{5}
\ifLuaTeX
  \usepackage{selnolig}  % disable illegal ligatures
\fi
\IfFileExists{bookmark.sty}{\usepackage{bookmark}}{\usepackage{hyperref}}
\IfFileExists{xurl.sty}{\usepackage{xurl}}{} % add URL line breaks if available
\urlstyle{same}
\hypersetup{
  colorlinks=true,
  linkcolor={blue},
  filecolor={Maroon},
  citecolor={Blue},
  urlcolor={blue},
  pdfcreator={LaTeX via pandoc}}

\author{}
\date{}

\begin{document}

\hypertarget{ux89e3ux7b54-10-ux767aux5c55ux30c8ux30d4ux30c3ux30af}{%
\section{解答 10 ---
発展トピック}\label{ux89e3ux7b54-10-ux767aux5c55ux30c8ux30d4ux30c3ux30af}}

\hypertarget{section}{%
\subsection{1}\label{section}}

固有値は1だけです。\texttt{(A-I)v=0} から固有空間は
\texttt{span\{(1,0)\^{}T\}}
の1次元しかなく、2本の独立固有ベクトルを持たないため対角化できません。

\hypertarget{section-1}{%
\subsection{2}\label{section-1}}

\texttt{A=I+N} で \texttt{I,N} は可換、\texttt{N\^{}2=0} なので

\[
e^{At}=e^t e^{Nt}=e^t(I+tN)
=e^t\begin{bmatrix}1&t\\0&1\end{bmatrix}.
\]

\hypertarget{section-2}{%
\subsection{3}\label{section-2}}

各固有方向は概ね \texttt{e\^{}\{λt\}}
で変化します。\texttt{Re(λ)\textless{}0}
なら減衰、\texttt{Re(λ)\textgreater{}0} なら増幅します。Jordan block
がある場合は多項式因子も現れます。

\hypertarget{section-3}{%
\subsection{4}\label{section-3}}

標準基底 \texttt{e\_i} に対する値 \texttt{a\_i=f(e\_i)}
を並べます。\texttt{x=Σx\_ie\_i} と線形性から
\texttt{f(x)=Σa\_ix\_i=a\^{}Tx} です。

\hypertarget{section-4}{%
\subsection{5}\label{section-4}}

\texttt{uv\^{}Tx=u(v\^{}Tx)} なので出力は常に \texttt{span\{u\}}
に入ります。よって像の次元は高々1です。SVD は一般行列をこの rank-1
行列の直交的な和として分解します。

\end{document}
